\section{Referencial Teórico}
\label{referencial}

% ORÇAMENTO: até 1 página. Template: conceitos com referência (de preferência mais de uma)
% + trabalhos correlatos, mostrando como as ideias evoluíram até o estado da arte.
% Citações em formato autor-ano (natbib): \citet{chave} -> Autor (ano); \citep{chave} -> (Autor, ano).
% Só citar o que alguém do grupo efetivamente leu. Conferir cada .bib (Google Acadêmico vem incompleto).

\subsection{Conceitos Fundamentais}
\label{ssec:conceitos}
% Só o que será usado depois, 1-2 frases cada:
% - grafo ponderado; caminho mínimo (Floyd-Warshall)
% - CENTRAL: problema dos carteiros. Carteiro Chinês (todas as arestas) e Carteiro Rural (só as obrigatórias);
%   grafo euleriano (grau par + conexidade) e por que arestas são repetidas
% - diferença para o caixeiro viajante (TSP): visitar vértices x cobrir arestas
% - algoritmo exato x heurística; aproximação (MST dobrada), se usada
% - coocorrência e medida normalizada (Jaccard)

\subsection{Trabalhos Correlatos}
\label{ssec:correlatos}
% Linha do tempo das ideias (ontogenia):
% Ratliff e Rosenthal (1983) -> Hall (1993) / Petersen (1997), heurísticas
% -> Roodbergen e de Koster (2001), várias travessas -> De Koster et al. (2007), revisão
% -> alocação de produtos correlacionados ("correlated storage assignment")
% Fechar dizendo onde ESTE trabalho se posiciona.
